\documentclass[11pt]{article}

\usepackage[margin=1in]{geometry}
\usepackage{amsmath, amssymb}
\usepackage{graphicx}
\usepackage{float}
\usepackage{booktabs}
\usepackage{xcolor}
\usepackage{listings}
\usepackage{hyperref}

\lstset{
  language=Python,
  basicstyle=\ttfamily\small,
  keywordstyle=\color{blue!70!black},
  commentstyle=\color{green!40!black},
  stringstyle=\color{red!60!black},
  frame=single,
  breaklines=true,
  showstringspaces=false,
  columns=fullflexible,
}

\title{MATH 7243 --- Lab 3: Linear Methods in Classification}
\author{Wilson Glass}
\date{October 21, 2026}

\begin{document}
\maketitle

%=============================================================================
\section*{Problem 1: Gender Recognition by Voice}

\subsection*{Setup}

The dataset contains 3168 voice recordings, each described by 20 acoustic features
(\texttt{meanfreq}, \texttt{sd}, \texttt{median}, \texttt{Q25}, \texttt{Q75},
\texttt{IQR}, \texttt{skew}, \texttt{kurt}, \texttt{sp.ent}, \texttt{sfm}, \texttt{mode},
\texttt{centroid}, \texttt{meanfun}, \texttt{minfun}, \texttt{maxfun}, \texttt{meandom},
\texttt{mindom}, \texttt{maxdom}, \texttt{dfrange}, \texttt{modindx}) and labelled
\texttt{male} or \texttt{female}. Following the Master notebook, every feature is
normalized by centering on its mean and dividing by its standard deviation.

\begin{lstlisting}
import pandas as pd
import numpy as np
from matplotlib import pyplot as plt
import seaborn as sns

data = pd.read_csv("voice.csv")

X = data.drop(columns=["label"])
y = data["label"]
X = (X - X.mean()) / X.std()      # normalize each feature
\end{lstlisting}

%-----------------------------------------------------------------------------
\subsection*{Question 1, Part 1: Exploratory Analysis}

\subsubsection*{Class balance}

\begin{lstlisting}
display(y.value_counts())
sns.countplot(x="label", data=data)
\end{lstlisting}

The classes are exactly balanced: 1584 male and 1584 female recordings. So a classifier
that always guesses one label gets 50\% accuracy, and plain accuracy is a fair metric
for this problem.

\subsubsection*{Correlation among the features}

\begin{lstlisting}
sns.heatmap(X.corr(), annot=True, linewidth=.5, fmt='.1f', annot_kws={"size": 7})
\end{lstlisting}

\begin{figure}[h]
  \centering
  \includegraphics[width=0.85\textwidth]{figures/p1_corr.png}
  \caption{Correlation between the 20 normalized features.}
\end{figure}

Many of the features are redundant:
\begin{itemize}
  \item \texttt{centroid} and \texttt{meanfreq} have correlation $1.0$. They are the same
  feature in this dataset.
  \item \texttt{skew} and \texttt{kurt} ($1.0$), and \texttt{maxdom} and \texttt{dfrange}
  ($1.0$), are also near-duplicates.
  \item The frequency-quantile features \texttt{meanfreq}, \texttt{median}, \texttt{Q25}
  and \texttt{centroid} form a block with correlations of $0.8$--$1.0$. \texttt{sd},
  \texttt{IQR}, \texttt{sp.ent} and \texttt{sfm} form a second block that is negatively
  correlated with the first.
  \item Several features are exact linear combinations of others:
  $\texttt{IQR} = \texttt{Q75} - \texttt{Q25}$ and
  $\texttt{dfrange} = \texttt{maxdom} - \texttt{mindom}$.
\end{itemize}
This redundancy matters for Question 2: adding all 20 features gives a model much less new
information than the count suggests.

\subsubsection*{Correlation within each class}

As in the Master notebook, we compute the correlation matrix separately for male and
female voices, on a shared color scale, and take their difference.

\begin{lstlisting}
I_m = y == "male"
I_f = y == "female"

vmin = min(X[I_m].corr().min().min(), X[I_f].corr().min().min())
vmax = max(X[I_m].corr().max().max(), X[I_f].corr().max().max())

f, axes = plt.subplots(1, 2, figsize=(20, 8))
sns.heatmap(X[I_m].corr(), annot=True, fmt='.1f', ax=axes[0], vmin=vmin, vmax=vmax)
sns.heatmap(X[I_f].corr(), annot=True, fmt='.1f', ax=axes[1], vmin=vmin, vmax=vmax)

sns.heatmap(X[I_m].corr() - X[I_f].corr(), annot=True, fmt='.1f')
\end{lstlisting}

\begin{figure}[h]
  \centering
  \includegraphics[width=\textwidth]{figures/p1_corr_by_class.png}
  \caption{Feature correlations for male (left) and female (right) voices.}
\end{figure}

\begin{figure}[h]
  \centering
  \includegraphics[width=0.7\textwidth]{figures/p1_corr_diff.png}
  \caption{Difference of the two class correlation matrices (male $-$ female).}
\end{figure}

Most of the correlation structure is the same for both genders: the difference matrix is
mostly within $\pm 0.2$. The largest differences involve \texttt{meanfun}, which is
correlated with \texttt{meandom}, \texttt{maxdom} and \texttt{dfrange} for male voices
(difference $0.4$--$0.5$) but not for female voices, and \texttt{sp.ent}, whose
correlation with \texttt{IQR} and \texttt{Q25} changes by up to $0.5$. Because the two
classes have different covariance structure, a method like QDA, which fits a separate
covariance matrix to each class, could have an advantage over methods that assume one
shared covariance.

\clearpage
\subsubsection*{Distributions by class}

\begin{lstlisting}
vio = pd.concat([y, X], axis=1)
vio = pd.melt(vio, id_vars="label", var_name="features", value_name="value")

sns.violinplot(x="features", y="value", hue="label", data=vio, split=True, inner="quart")
sns.boxplot(x="features", y="value", hue="label", data=vio)
\end{lstlisting}

\begin{figure}[h]
  \centering
  \includegraphics[width=\textwidth]{figures/p1_violin_box_0.png}
  \includegraphics[width=\textwidth]{figures/p1_violin_box_1.png}
  \caption{Violin plots (top) and box plots (bottom) of each normalized feature, split by
  label.}
\end{figure}

The Master notebook's swarm plot did not work for this data. Several features
(\texttt{maxfun}, \texttt{mindom}, \texttt{modindx}) take only a few distinct values, and a
swarm plot cannot place that many tied points side by side: seaborn reported that up to
70\% of the points could not be placed. We use a jittered strip plot on a random subsample
of 500 recordings instead.

\begin{lstlisting}
vio_s = pd.concat([y, X], axis=1).sample(500, random_state=0)
vio_s = pd.melt(vio_s, id_vars="label", var_name="features", value_name="value")

sns.stripplot(x="features", y="value", hue="label", data=vio_s,
              size=2, jitter=0.3, dodge=True, alpha=.6)
\end{lstlisting}

\begin{figure}[h]
  \centering
  \includegraphics[width=\textwidth]{figures/p1_strip.png}
  \caption{Strip plot of 500 randomly chosen recordings.}
\end{figure}

\texttt{meanfun} stands out in all three plots: the male and female distributions barely
overlap, with male voices well below the mean and female voices above it. \texttt{IQR},
\texttt{Q25}, \texttt{sd}, \texttt{sp.ent} and \texttt{sfm} also show clearly shifted
medians, but with much more overlap. Features such as \texttt{skew}, \texttt{kurt},
\texttt{Q75} and \texttt{modindx} look almost the same for both classes.

%-----------------------------------------------------------------------------
\clearpage
\subsection*{Question 1, Part 2: Most Indicative Features}

To rank the features by how well they separate the genders, we compute three measures for
each one:
\begin{itemize}
  \item the gap between the male and female means (in standard deviations, since $X$ is
  normalized);
  \item the absolute correlation with the 0/1 label;
  \item the 5-fold cross-validated accuracy of a logistic regression that uses only that
  one feature.
\end{itemize}

\begin{lstlisting}
from sklearn.linear_model import LogisticRegression
from sklearn.model_selection import cross_val_score

y_num = (y == "male").astype(int)

rank = pd.DataFrame({
    "mean_diff": (X[I_m].mean() - X[I_f].mean()).abs(),
    "corr_w_label": X.corrwith(y_num).abs(),
    "single_feat_acc": [cross_val_score(LogisticRegression(), X[[c]], y, cv=5).mean()
                        for c in X.columns],
}).sort_values("single_feat_acc", ascending=False)
\end{lstlisting}

\begin{center}
\begin{tabular}{lccc}
\toprule
Feature & Mean gap (SD) & $|$corr with label$|$ & Single-feature CV accuracy \\
\midrule
\texttt{meanfun} & 1.668 & 0.834 & 0.950 \\
\texttt{IQR}     & 1.238 & 0.619 & 0.887 \\
\texttt{Q25}     & 1.023 & 0.511 & 0.821 \\
\texttt{sd}      & 0.959 & 0.480 & 0.754 \\
\texttt{sp.ent}  & 0.981 & 0.491 & 0.735 \\
\texttt{sfm}     & 0.715 & 0.357 & 0.651 \\
\midrule
\texttt{skew}    & 0.073 & 0.037 & 0.448 \\
\texttt{modindx} & 0.062 & 0.031 & 0.441 \\
\bottomrule
\end{tabular}
\end{center}

\begin{figure}[h]
  \centering
  \includegraphics[width=0.5\textwidth]{figures/p1_top2.png}
  \caption{The two most indicative features, \texttt{meanfun} and \texttt{IQR}.}
\end{figure}

\textbf{Answer:} The two most indicative features are \textbf{\texttt{meanfun}} and
\textbf{\texttt{IQR}}. All three measures rank them first and second.
\begin{itemize}
  \item \texttt{meanfun}, the average fundamental frequency, measures the pitch of the
  voice. On its own it classifies 95.0\% of recordings correctly, and its class means are
  1.67 standard deviations apart.
  \item \texttt{IQR}, the interquartile range of the frequency spectrum, gets 88.7\%
  alone. Male voices have a wider spread of frequencies (higher \texttt{IQR}).
  \item \texttt{Q25} is a close third, but it is strongly correlated with \texttt{IQR}
  ($-0.9$), so it would add little as a second feature. \texttt{IQR} is only moderately
  correlated with \texttt{meanfun} ($-0.5$), so it does add new information.
\end{itemize}
In the scatter plot the two classes form two clearly separated clusters: males at low
\texttt{meanfun} and high \texttt{IQR}, females at high \texttt{meanfun} and low
\texttt{IQR}.

%-----------------------------------------------------------------------------
\clearpage
\subsection*{Question 2: Linear Regression, Logistic Regression and QDA}

\subsubsection*{Setup}

We split the data once into 70\% training and 30\% test data (stratified by label), so that
the two-feature and all-feature models are evaluated on the same 951 test recordings.

Linear regression is used as a classifier as in the Master notebook. We regress on the
one-hot (dummy) labels and predict the class whose column has the larger fitted value. We
wrap it in a small class so it has the same \texttt{predict}/\texttt{score} interface as
the sklearn classifiers.

\begin{lstlisting}
from sklearn.linear_model import LinearRegression, LogisticRegression
from sklearn.discriminant_analysis import QuadraticDiscriminantAnalysis
from sklearn.model_selection import train_test_split

X_train, X_test, y_train, y_test = train_test_split(X, y, test_size=0.3,
                                                    stratify=y, random_state=0)
feats2 = ["meanfun", "IQR"]

class LinRegClassifier:
    def fit(self, X, y):
        Y = pd.get_dummies(y)
        self.classes_ = Y.columns.to_numpy()
        self.lr = LinearRegression().fit(X, Y.astype(float))
        return self

    def predict(self, X):
        return self.classes_[np.argmax(self.lr.predict(X), axis=1)]

    def predict_proba(self, X):
        return self.lr.predict(X)

    def score(self, X, y):
        return np.mean(self.predict(X) == np.asarray(y))
\end{lstlisting}

\subsubsection*{Two-feature fits}

\begin{lstlisting}
models2 = {
    "Linear Regression": LinRegClassifier(),
    "Logistic Regression": LogisticRegression(),
    "QDA": QuadraticDiscriminantAnalysis(store_covariance=True),
}
for name, m in models2.items():
    m.fit(X_train[feats2], y_train)
\end{lstlisting}

The decision regions are drawn by predicting on a $200 \times 200$ grid covering the data
and shading each grid point by its predicted label.

\begin{figure}[h]
  \centering
  \includegraphics[width=\textwidth]{figures/p1_boundaries.png}
  \caption{Decision boundaries on \texttt{meanfun} and \texttt{IQR}, with the test
  recordings. Blue region: predicted male; orange region: predicted female.}
\end{figure}

Linear and logistic regression draw almost the same straight line. It tilts upward: a voice
with a larger \texttt{IQR} needs a higher \texttt{meanfun} to be labelled female. QDA's
boundary is a curve that wraps around the male cluster. The male cluster is much tighter
than the female one, and QDA's separate class covariances let the boundary bend around it.
A small second group of females at high \texttt{IQR} (around 2) is on the female side
for every model. Most errors come from the region where the clusters overlap, near
$\texttt{meanfun} \approx 0$.

\subsubsection*{All-feature fits}

Since \texttt{centroid} $=$ \texttt{meanfreq}, $\texttt{IQR} = \texttt{Q75} - \texttt{Q25}$
and $\texttt{dfrange} = \texttt{maxdom} - \texttt{mindom}$, the class covariance matrices
are singular when all 20 features are used. A small regularization parameter keeps QDA well
defined, and logistic regression is given more iterations so it converges.

\begin{lstlisting}
models_all = {
    "Linear Regression": LinRegClassifier(),
    "Logistic Regression": LogisticRegression(max_iter=1000),
    "QDA": QuadraticDiscriminantAnalysis(reg_param=0.01),
}
for name, m in models_all.items():
    m.fit(X_train, y_train)
\end{lstlisting}

A 20-dimensional decision boundary cannot be drawn. Instead, we graph each model's fitted
``male'' score on the test set (the predicted probability for logistic regression and QDA,
the fitted dummy value for linear regression), split by true label.

\begin{figure}[h]
  \centering
  \includegraphics[width=\textwidth]{figures/p1_score_hist.png}
  \caption{Fitted male score on the test set by true label, for two features (top) and all
  features (bottom). The dashed line is the decision threshold.}
\end{figure}

\subsubsection*{Results}

\begin{center}
\begin{tabular}{lcccc}
\toprule
 & \multicolumn{2}{c}{2 features} & \multicolumn{2}{c}{All 20 features} \\
\cmidrule(lr){2-3} \cmidrule(lr){4-5}
Model & Train & Test & Train & Test \\
\midrule
Linear Regression   & 0.965 & 0.956 & 0.973 & 0.962 \\
Logistic Regression & 0.968 & 0.960 & 0.978 & \textbf{0.965} \\
QDA                 & 0.969 & \textbf{0.962} & 0.973 & 0.959 \\
\bottomrule
\end{tabular}
\end{center}

\begin{center}
\begin{tabular}{lcccc}
\toprule
 & \multicolumn{2}{c}{2 features} & \multicolumn{2}{c}{All 20 features} \\
\cmidrule(lr){2-3} \cmidrule(lr){4-5}
Model & male $\to$ female & female $\to$ male & male $\to$ female & female $\to$ male \\
\midrule
Linear Regression   & 13 & 29 & 15 & 21 \\
Logistic Regression & 16 & 22 & 17 & 16 \\
QDA                 & 17 & 19 & 22 & 17 \\
\bottomrule
\end{tabular}
\end{center}
\noindent (Number of the 951 test recordings misclassified in each direction.)

\subsubsection*{How the two-feature fit compares to the fit on all features}

\begin{itemize}
  \item \textbf{Two features already give almost all of the accuracy.} With only
  \texttt{meanfun} and \texttt{IQR}, every method classifies about 96\% of test recordings
  correctly. QDA is the best two-feature model (96.2\%), but only slightly.
  \item \textbf{All 20 features add very little.} Linear and logistic regression improve by
  only about half a percentage point (to 96.2\% and 96.5\%), and QDA gets slightly
  \emph{worse} (96.2\% $\to$ 95.9\%). As the correlation heatmap shows, the other features
  are mostly near-copies of each other or carry little class information, so they add almost
  nothing beyond \texttt{meanfun}.
  \item \textbf{QDA overfits with all features.} It has to estimate two full $20 \times 20$
  covariance matrices. Its train/test gap grows from 0.7 points to 1.4 points. Its score
  histogram shows the same thing: it pushes nearly every recording to a probability of
  exactly 0 or 1, including the ones it gets wrong.
  \item \textbf{Logistic regression benefits most} from the extra features, cutting females
  mislabelled as male from 22 to 16. Linear regression's errors are lopsided with two
  features (29 females labelled male vs.\ 13 males labelled female), and the extra features
  partly correct this.
  \item Linear regression's scores are not probabilities. They run from about $-0.6$ to
  $1.6$, and its overlap around the threshold is wider than logistic regression's.
  \item \textbf{Caveat:} with 951 test points, one standard error of accuracy is about
  $\sqrt{0.04 \cdot 0.96 / 951} \approx 0.6$ points. So the differences between the three
  methods, and between 2 and 20 features, are all within noise for this single split.
\end{itemize}

\textbf{Conclusion.} Two well-chosen features give essentially all of the classification
power (about 96\%), and a model whose decision boundary can be plotted and interpreted.
Adding the other 18 features buys at most about half a point of accuracy, makes the model
uninterpretable, and in the case of QDA introduces overfitting.

%=============================================================================
\clearpage
\section*{Problem 2: Alzheimer MRI Classification}

\subsection*{Setup}

The Alzheimer4 dataset contains MRI brain slices in four classes, already split into
training and test folders:

\begin{center}
\begin{tabular}{lcccc}
\toprule
 & NonDemented & VeryMildDemented & MildDemented & ModerateDemented \\
\midrule
Label & 0 & 1 & 2 & 3 \\
Train & 2560 & 1792 & 717 & 52 \\
Test  & 640  & 448  & 179 & 12 \\
\bottomrule
\end{tabular}
\end{center}

The classes are very unbalanced: NonDemented is half of the data and ModerateDemented
only 1\%. A model that always predicts NonDemented gets 50\% accuracy. So besides accuracy we
report \textbf{balanced accuracy}, the average of the four per-class recalls, for which
that trivial model scores 25\%.

Each image is grayscale, $176 \times 208$ pixels, and is resized to $128 \times 128$. The
lab's starter code repeated the same loading loop for each of the four class folders. We
replace it with a single function that loads a whole split.

\begin{lstlisting}
DATA_DIR = "data/AlzheimerDataset4/Alzheimer_s Dataset"
CLASSES = ["NonDemented", "VeryMildDemented", "MildDemented", "ModerateDemented"]
IMG_SIZE = (128, 128)

def load_split(split):
    data, labels, paths = [], [], []
    for label, cls in enumerate(CLASSES):
        folder = os.path.join(DATA_DIR, split, cls)
        for filename in sorted(os.listdir(folder)):
            path = os.path.join(folder, filename)
            img = np.array(Image.open(path).resize(IMG_SIZE))
            if img.shape == IMG_SIZE:          # skip any non-grayscale images
                data.append(img)
                labels.append(label)
                paths.append(path)
    return np.array(data), np.array(labels), paths

train_data, train_y, train_paths = load_split("train")
train_label = np.eye(len(CLASSES))[train_y]          # one-hot, shape (5121, 4)
\end{lstlisting}

%-----------------------------------------------------------------------------
\subsection*{Question 1: Import the Test Dataset}

With the loader function, importing the test set is a single call:

\begin{lstlisting}
test_data, test_y, test_paths = load_split("test")
test_label = np.eye(len(CLASSES))[test_y]
\end{lstlisting}

This gives \texttt{test\_data} of shape $(1279, 128, 128)$ and \texttt{test\_label} of
shape $(1279, 4)$, matching the expected size. Training and test images come from the
provided folders, so the split was fixed before any processing.

\begin{figure}[h]
  \centering
  \includegraphics[width=0.45\textwidth]{figures/p2_example.png}
  \caption{One training image. Pixel values range from 0 to 255.}
\end{figure}

%-----------------------------------------------------------------------------
\clearpage
\subsection*{Question 2: LDA, QDA, Logistic Regression and Categorical Linear Regression}

\subsubsection*{Setup}

Each $128 \times 128$ image is flattened into a row of 16{,}384 pixels and scaled to
$[0, 1]$. Categorical linear regression regresses on the one-hot labels and predicts the
class with the largest fitted value.

\begin{lstlisting}
from sklearn.discriminant_analysis import LinearDiscriminantAnalysis, QuadraticDiscriminantAnalysis
from sklearn.linear_model import LinearRegression, LogisticRegression
from sklearn.decomposition import PCA
from sklearn.pipeline import make_pipeline

X_train = (train_data.reshape(len(train_data), -1) / 255).astype(np.float32)
X_test  = (test_data.reshape(len(test_data), -1) / 255).astype(np.float32)

class LinRegClassifier:
    def fit(self, X, y):
        self.lr = LinearRegression().fit(X, np.eye(len(CLASSES))[y])
        return self

    def predict(self, X):
        return np.argmax(self.lr.predict(X), axis=1)

models = {
    "LDA": LinearDiscriminantAnalysis(),
    "QDA (on 40 PCA components)": make_pipeline(PCA(n_components=40, random_state=0),
                                                QuadraticDiscriminantAnalysis(reg_param=0.1)),
    "Logistic Regression": LogisticRegression(max_iter=500),
    "Categorical Linear Regression": LinRegClassifier(),
}
for name, m in models.items():
    m.fit(X_train, train_y)
\end{lstlisting}

\textbf{QDA needs dimension reduction.} QDA estimates a separate $16384 \times 16384$
covariance matrix for each class, but every class has fewer images than pixels, so these
matrices are singular. scikit-learn refuses to fit it (\texttt{LinAlgError: The covariance
matrix of class 0 is not full rank}). Its suggested alternative, the \texttt{eigen} solver
with shrinkage, would have to decompose four $16384 \times 16384$ matrices, which is
impractical. Instead, we project the images onto their top 40 principal components (PCA fitted
on the training set only) and run QDA there. We chose 40 because the smallest class,
ModerateDemented, has just 52 training images, and each class needs more samples than
dimensions.

\subsubsection*{Results}

\begin{center}
\begin{tabular}{lcccc}
\toprule
Model & Train acc & Test acc & Train balanced & Test balanced \\
\midrule
LDA                           & 1.000 & 0.396 & 1.000 & 0.290 \\
QDA (40 PCA components)       & 0.793 & 0.535 & 0.851 & 0.369 \\
Logistic Regression           & 1.000 & \textbf{0.598} & 1.000 & 0.365 \\
Categorical Linear Regression & 1.000 & 0.561 & 1.000 & \textbf{0.449} \\
\midrule
Always ``NonDemented''        & --- & 0.500 & --- & 0.250 \\
\bottomrule
\end{tabular}
\end{center}

\begin{figure}[p]
  \centering
  \includegraphics[width=0.82\textwidth]{figures/p2_confusion.png}
  \caption{Confusion matrices on the training set (left) and test set (right). Rows are the
  true class, columns the predicted class.}
\end{figure}

\subsubsection*{Confusion matrices}

Rows are the true class and columns the predicted class, in the order Non, VeryMild, Mild,
Moderate.
\[
\begin{array}{lcc}
 & \text{Train} & \text{Test} \\[4pt]
\text{LDA} &
\begin{pmatrix} 2560&0&0&0\\0&1792&0&0\\0&0&717&0\\0&0&0&52 \end{pmatrix} &
\begin{pmatrix} 480&12&15&133\\304&14&17&113\\108&6&8&57\\7&1&0&4 \end{pmatrix} \\[24pt]
\text{QDA} &
\begin{pmatrix} 2115&318&127&0\\423&1275&94&0\\36&60&621&0\\0&0&0&52 \end{pmatrix} &
\begin{pmatrix} 411&221&8&0\\192&252&4&0\\46&114&19&0\\0&10&0&2 \end{pmatrix} \\[24pt]
\text{Logistic} &
\begin{pmatrix} 2560&0&0&0\\0&1792&0&0\\0&0&717&0\\0&0&0&52 \end{pmatrix} &
\begin{pmatrix} 497&140&3&0\\202&242&4&0\\71&82&26&0\\6&6&0&0 \end{pmatrix} \\[24pt]
\text{Linear Reg.} &
\begin{pmatrix} 2560&0&0&0\\0&1792&0&0\\0&0&717&0\\0&0&0&52 \end{pmatrix} &
\begin{pmatrix} 425&197&18&0\\186&242&20&0\\59&74&46&0\\6&2&0&4 \end{pmatrix}
\end{array}
\]

The per-class test recall (the diagonal divided by the row total) is:

\begin{center}
\begin{tabular}{lcccc}
\toprule
Model & Non & VeryMild & Mild & Moderate \\
\midrule
LDA                 & 0.75 & 0.03 & 0.04 & 0.33 \\
QDA                 & 0.64 & 0.56 & 0.11 & 0.17 \\
Logistic Regression & 0.78 & 0.54 & 0.15 & 0.00 \\
Linear Regression   & 0.66 & 0.54 & 0.26 & 0.33 \\
\bottomrule
\end{tabular}
\end{center}

\subsubsection*{Discussion}

\begin{itemize}
  \item \textbf{Overfitting dominates.} There are 16{,}384 pixel features but only 5121
  training images. LDA, logistic regression and linear regression all reach 100\% training
  accuracy, with perfectly diagonal training confusion matrices, but only 40--60\% on the test
  set. With more features than samples the training classes are linearly separable, so these
  models memorize the training images instead of learning patterns that generalize. QDA on 40
  PCA components has far fewer effective parameters. It is the only model that does not fit
  the training set perfectly (79\%), and it has the smallest train/test gap (26 points vs.\
  40--60).
  \item \textbf{Test performance is only modestly better than guessing.} The best test
  accuracy, logistic regression's 59.8\%, is about 10 points above the 50\% baseline. LDA does
  \emph{worse} than the baseline: it labels 307 test images as ModerateDemented, although
  only 12 are. Its discriminant directions were fitted to noise in the 16{,}384-dimensional
  pixel space and do not transfer to new images.
  \item \textbf{Accuracy and balanced accuracy rank the models differently.} Logistic
  regression has the best accuracy because it favours the large classes, but it never
  predicts ModerateDemented (0 of 12 correct). Categorical linear regression has the best
  balanced accuracy (0.449). It recovers the most MildDemented images (26\%) and 4 of the 12
  ModerateDemented images, at the cost of more NonDemented errors.
  \item \textbf{Errors are mostly between neighbouring stages.} In every test confusion
  matrix, most mistakes are NonDemented $\leftrightarrow$ VeryMildDemented, and MildDemented
  images are mostly labelled VeryMild or Non rather than Moderate. Dementia severity is a
  continuum, and adjacent stages look very similar on a single MRI slice. The models largely
  separate ``no'' from ``some'' dementia but struggle to grade the severity.
\end{itemize}

\textbf{Conclusion.} No linear method works well on raw pixels. Logistic regression gives
the highest accuracy and categorical linear regression the best balance across classes. QDA
is only feasible after dimension reduction, which also makes it the least overfit. LDA is
the worst, since it cannot estimate a pooled covariance from 5121 samples in 16{,}384
dimensions. Reducing the number of features (downsampling or PCA), regularizing, or
reweighting the classes (\texttt{class\_weight="balanced"}) would be natural next steps.

%-----------------------------------------------------------------------------
\clearpage
\subsection*{Question 3 (Bonus): Downsampling}

\subsubsection*{Modifying the downsample function}

The provided \texttt{downsample} function splits each image into \texttt{DS} new images by
taking every \texttt{DS}-th pixel with a different starting offset:
row $n \cdot \texttt{DS} + i$ of the new array is \texttt{X[n, i::DS]}. We made four changes:
\begin{itemize}
  \item The image size is read from \texttt{X} instead of being hard-coded as 123{,}008.
  Our flattened images have $128 \times 128 = 16{,}384$ pixels.
  \item \texttt{y} is a numpy array here, so \texttt{y.to\_list()} is no longer needed.
  \item The double loop is replaced by an equivalent reshape, which is much faster.
  \item The function also returns which original image each new row came from, so that the
  \texttt{DS} predictions for each test image can be combined by majority vote.
\end{itemize}

\begin{lstlisting}
def downsample(X, y, DS):
    N, im_size = X.shape

    if im_size % DS > 0:
        print("Downsample rate is not a divisor of", im_size)
        DS = 1

    # (N, im_size) -> (N, im_size/DS, DS) -> swap last two axes -> rows X[n, i::DS]
    Xds = X.reshape(N, im_size // DS, DS).transpose(0, 2, 1).reshape(N * DS, im_size // DS)
    yds = np.repeat(y, DS)
    groups = np.repeat(np.arange(N), DS)

    return Xds, yds, groups

DS = 8
Xds_train, yds_train, g_train = downsample(X_train, train_y, DS)
Xds_test,  yds_test,  g_test  = downsample(X_test,  test_y,  DS)

assert np.array_equal(Xds_train[3*DS + 5], X_train[3, 5::DS])   # matches the loop definition
\end{lstlisting}

With $\texttt{DS} = 8$ each image becomes 8 images of 2048 pixels, giving 40{,}968 training
rows and 10{,}232 test rows. As the lab requires, train and test were split
\emph{before} downsampling, so no sub-image of a test image is ever seen in training. Since
128 is divisible by 8, taking every 8th pixel of the flattened image keeps every 8th
\emph{column} in each row. So each new image is $128 \times 16$ pixels, and the 8 images
are horizontally shifted copies of the original.

\begin{figure}[H]
  \centering
  \vspace{0.5em}
  \includegraphics[width=0.7\textwidth]{figures/p2_ds_example.png}
  \caption{An original image and two of its eight downsampled versions.}
\end{figure}

\clearpage
\subsubsection*{Models on the downsampled data}

The same four models are fit to the downsampled rows. QDA still uses 40 PCA components,
because the smallest class has only $52 \times 8 = 416$ rows, fewer than the 2048 features.
Each test image is then classified by a majority vote over its 8 sub-images.

\begin{lstlisting}
def vote(pred, groups):
    counts = np.zeros((groups.max() + 1, len(CLASSES)))
    np.add.at(counts, (groups, pred.astype(int)), 1)
    return counts.argmax(axis=1)

for name, m in models_ds.items():
    m.fit(Xds_train, yds_train)
    p_train, p_test = m.predict(Xds_train), m.predict(Xds_test)
    img_train, img_test = vote(p_train, g_train), vote(p_test, g_test)
\end{lstlisting}

\subsubsection*{Results}

\begin{center}
\begin{tabular}{lccccc}
\toprule
 & \multicolumn{2}{c}{Sub-image accuracy} & \multicolumn{3}{c}{Image level (vote)} \\
\cmidrule(lr){2-3} \cmidrule(lr){4-6}
Model & Train & Test & Train acc & Test acc & Test balanced \\
\midrule
LDA                           & 0.783 & 0.517 & 0.862 & 0.540 & 0.326 \\
QDA (40 PCA components)       & 0.642 & 0.551 & 0.664 & 0.543 & \textbf{0.392} \\
Logistic Regression           & 0.778 & 0.521 & 0.855 & \textbf{0.558} & 0.323 \\
Categorical Linear Regression & 0.769 & 0.519 & 0.834 & 0.532 & 0.296 \\
\bottomrule
\end{tabular}
\end{center}

\begin{center}
\begin{tabular}{lcccc}
\toprule
 & \multicolumn{2}{c}{Test accuracy} & \multicolumn{2}{c}{Test balanced accuracy} \\
\cmidrule(lr){2-3} \cmidrule(lr){4-5}
Model & Full resolution (Q2) & $\texttt{DS} = 8$ & Full resolution (Q2) & $\texttt{DS} = 8$ \\
\midrule
LDA                           & 0.396 & 0.540 & 0.290 & 0.326 \\
QDA (40 PCA components)       & 0.535 & 0.543 & 0.369 & 0.392 \\
Logistic Regression           & 0.598 & 0.558 & 0.365 & 0.323 \\
Categorical Linear Regression & 0.561 & 0.532 & 0.449 & 0.296 \\
\bottomrule
\end{tabular}
\end{center}

\begin{figure}[p]
  \centering
  \includegraphics[width=0.82\textwidth]{figures/p2_ds_confusion.png}
  \caption{Confusion matrices after downsampling ($\texttt{DS} = 8$) and majority voting,
  on the training set (left) and test set (right).}
\end{figure}

\[
\begin{array}{lcc}
 & \text{Train (vote)} & \text{Test (vote)} \\[4pt]
\text{LDA} &
\begin{pmatrix} 2423&124&13&0\\404&1374&14&0\\64&86&567&0\\0&0&0&52 \end{pmatrix} &
\begin{pmatrix} 504&128&8&0\\257&181&10&0\\73&101&5&0\\4&7&0&1 \end{pmatrix} \\[24pt]
\text{QDA} &
\begin{pmatrix} 2035&320&201&4\\746&866&177&3\\136&128&448&5\\0&0&0&52 \end{pmatrix} &
\begin{pmatrix} 516&124&0&0\\282&164&2&0\\67&101&11&0\\3&5&0&4 \end{pmatrix} \\[24pt]
\text{Logistic} &
\begin{pmatrix} 2408&133&19&0\\410&1366&16&0\\74&89&554&0\\0&0&0&52 \end{pmatrix} &
\begin{pmatrix} 513&121&6&0\\242&188&18&0\\70&96&13&0\\6&6&0&0 \end{pmatrix} \\[24pt]
\text{Linear Reg.} &
\begin{pmatrix} 2466&90&4&0\\458&1332&2&0\\116&142&459&0\\10&26&3&13 \end{pmatrix} &
\begin{pmatrix} 504&132&4&0\\269&175&4&0\\80&98&1&0\\6&6&0&0 \end{pmatrix}
\end{array}
\]

\begin{center}
\begin{tabular}{lcccc}
\toprule
Model ($\texttt{DS} = 8$) & Non & VeryMild & Mild & Moderate \\
\midrule
LDA                 & 0.79 & 0.40 & 0.03 & 0.08 \\
QDA                 & 0.81 & 0.37 & 0.06 & 0.33 \\
Logistic Regression & 0.80 & 0.42 & 0.07 & 0.00 \\
Linear Regression   & 0.79 & 0.39 & 0.01 & 0.00 \\
\bottomrule
\end{tabular}
\end{center}
\noindent (Per-class test recall after voting.)

\subsubsection*{Discussion}

\begin{itemize}
  \item \textbf{Downsampling fixes the more-features-than-samples problem.} With 40{,}968 rows
  and 2048 features, no model fits the training set perfectly any more: image-level train
  accuracy is 66--86\% instead of 100\%. The sub-image train/test gap shrinks from 40--60
  points in Question 2 to about 25 points for LDA, logistic and linear regression, and 9
  points for QDA. LDA benefits most. Its pooled covariance can now be estimated, it no longer
  floods the predictions with ModerateDemented, and its test accuracy rises from 39.6\% to
  54.0\%.
  \item \textbf{But generalization does not improve overall.} All four models now reach
  53--56\% test accuracy, only a few points above the 50\% baseline. The best model gets
  worse: logistic regression drops from 59.8\% to 55.8\%. Balanced accuracy falls for
  logistic regression and sharply for linear regression (0.449 $\to$ 0.296). The models now
  almost never recognize MildDemented (at most 7\% recall, vs.\ up to 26\% before) and
  predict nearly everything as NonDemented or VeryMildDemented. The problem has moved from
  overfitting (variance) to lost information (bias).
  \item \textbf{Why.} The 8 sub-images of one MRI are nearly identical shifted copies, so the
  dataset is not really 8 times larger. It still holds information from only 5121 distinct
  images, and only 52 ModerateDemented ones. What changes is that each sub-image keeps only
  every 8th column, discarding most of the horizontal detail that separates neighbouring
  stages. The fewer features act as regularization, which helps the unstable LDA, but they
  also remove the signal the better models were using. The copies of the rare class are still
  memorized: LDA, QDA and logistic regression classify all 52 ModerateDemented training images
  correctly but get only 0--4 of the 12 test images right.
  \item \textbf{Majority voting helps a little.} Combining the 8 sub-image predictions raises
  test accuracy by 1--4 points for LDA, logistic and linear regression (e.g.\ logistic
  52.1\% $\to$ 55.8\%), and slightly lowers it for QDA (55.1\% $\to$ 54.3\%). Errors on the
  8 copies are strongly correlated, so voting cannot fix much.
  \item \textbf{Caveats.} Logistic regression reached \texttt{max\_iter = 500} without
  converging, so its numbers could change slightly with more iterations. All results come
  from a single train/test split.
\end{itemize}

\textbf{Conclusion.} Downsampling with $\texttt{DS} = 8$ makes every method numerically well
posed and removes the extreme overfitting of Question 2, which matters most for LDA. It
does not make the models better at recognizing dementia: the ``new'' images carry no new
information, and keeping every 8th column discards useful detail. Better options would be
to downsample in both directions (for example, averaging $2 \times 2$ or $4 \times 4$ pixel
blocks, which keeps the image shape), to try a smaller \texttt{DS}, to use PCA, or to
reweight the classes so the rare dementia stages are not ignored.

\end{document}
